\documentclass[12pt]{article}
\usepackage[utf8]{inputenc}
\usepackage[T1]{fontenc}
\usepackage[margin=1in]{geometry}
\usepackage{setspace}
\usepackage{xcolor}
\usepackage[colorlinks=true, urlcolor=blue!60!black, linkcolor=blue!60!black]{hyperref}
\setlength{\parindent}{0em}
\setlength{\parskip}{0.45em}
\renewcommand{\baselinestretch}{1.0}
\begin{document}
\begin{center}
{\Large \textit{Craig v. Missouri}, 29 U.S. (4 Pet.) 410 (1830)\par}
\vspace{0.3em}
{\normalsize Case Brief --- One-Page Edition\par}
\end{center}

\textbf{Memory Jogger.} State loan-office certificates as forbidden ``bills of credit'' (\href{https://supreme.justia.com/cases/federal/us/29/410/}{29 U.S. 410}).

\textbf{Facts.} After the Panic of 1819 and amid a shortage of specie, Missouri authorized up to \$200,000 in certificates of fifty cents to ten dollars, bearing 2\% interest, lent to citizens, receivable for taxes, salaries, and salt, and backed by a redemption fund and ``the faith of the state.'' They were not legal tender but were designed to circulate. Hiram Craig and two others gave the State a \$199.99 note for certificates and defaulted.

\textbf{Procedural History.} Missouri's courts enforced the note; writ of error.

\textbf{Votes.} Reversed, 4--3. Marshall wrote for the Court; Johnson, Thompson, and McLean dissented.

\textbf{Holding.} State paper issued on the state's faith to circulate as money is a bill of credit even without tender status, and a note given for it is void.

\textbf{Reasoning.} The clause targets state paper currency, not ordinary borrowing; the Constitution prohibits things, not names; and the separate tender clause shows bills of credit need not be tender. The dissenters stressed the interest, the pledged fund, and federalism. \textit{Briscoe} (1837) later upheld notes of a state-owned bank.

\textbf{Examples.} \emph{Same-Side:} faith-backed state certificates receivable for taxes. \emph{Opposite-Side:} ordinary state bonds. \emph{Fence-Sitter:} a reserve-backed state digital token.

\textbf{Critique.} \textit{Briscoe}'s corporate-form loophole undercut \textit{Craig}'s substance-over-form principle, and historians see Missouri's capped certificates as a modest response to specie scarcity.

\textbf{Key Quotation.} ``To emit bills of credit conveys to the mind the idea of issuing paper, intended to circulate through the community for its ordinary purposes as money, which paper is redeemable at a future day; this is the sense in which the terms have been always understood.'' --- Chief Justice Marshall, Opinion of the Court, 29 U.S. at 432.

\end{document}
